\section{Robustness, Scope, and Verification}\label{sec:robustness}

The quadratic-uniform baseline is chosen to expose the architecture-state feedback in closed form. The current theorem package is intentionally narrower than the set of numerical exercises retained in the repository.

\subsection{What is and is not generalized}

Propositions~\ref{prop:T1}--\ref{prop:T3} are quadratic-baseline theorems on strict $\mathcal R^+$. The abstract order implication
\[
 h_M<h_F,\quad \Psi'(h)>0
 \quad\Longrightarrow\quad
 \Psi(h_M)<\Psi(h_F)
\]
is useful for organizing the mechanism, but it is not a primitive demand theorem. In particular, strict concavity of future utility alone does not guarantee the state ordering $h_M<h_F$. A retained counterexample from the hostile generality audit reverses that ordering. We therefore make no arbitrary-concave-demand claim.

Likewise, the pre-repair numerical checks for alternative integration-cost CDFs are not carried into the current certified theorem set. No non-uniform-CDF theorem is asserted in this paper.

\subsection{Nonquadratic numerical robustness}

The repository retains deterministic numerical checks for common-curvature power specifications
\begin{equation}
 v_j(q)=a_jq-\frac{q^r}{r},
 \qquad r\in\{3,4\}.
 \label{eq:power-utility}
\end{equation}
For the tested parameterizations, anticipated metering continues to produce a lower high-use installed state and the computed metering-gain map remains increasing over the tested interval. These exercises are \emph{numerical robustness only}. They neither establish a theorem for arbitrary $r$ nor expand the quantifiers of Propositions~\ref{prop:T1}--\ref{prop:T3}.

\subsection{Why strict $\mathcal R^+$ is substantive}

The provider's continuation problem is globally defined at every installed state by equations \eqref{eq:VF-global} and \eqref{eq:VM-global}, but the closed-form headline theorem requires both-served continuation to dominate over the complete candidate interval $[h_0,h_F]$. This is exactly what strict $\mathcal R^+$ certifies.

Outside that domain, the active set can change. A permanent exact counterexample fixes
\begin{equation}
 a_L=4,\quad a_H=5,\quad c=1,\quad l=0.1,\quad h=0.6.
 \label{eq:scope-counterexample}
\end{equation}
At this installed state, the provider's $H$-only flat continuation exceeds the both-served flat continuation by exactly $23/10$. Hence a theorem quantified over all positive installed compositions is false. The counterexample is retained as a regression guard; it is not used to characterize the full equilibrium correspondence outside $\mathcal R^+$.

\subsection{Formal-verification boundary}

A targeted Lean 4 artifact certifies a proof-critical logical and algebraic core, including threshold ordering from state ordering and monotonicity, fixed-point uniqueness skeletons, strict-gap response reversal, positivity of the baseline derivative expression, repaired state-ordering components, the fixed-installed-base welfare identity, and the exact $23/10$ arithmetic scope guard.

The formal artifact does not encode the complete economic model. In particular, weak participation at zero surplus, the provider's complete active-set optimization, the atomless integration game, the full existence proof for the mixed equilibrium, endogenous welfare comparisons, and institutional interpretation are not fully formalized. Those claims are supported by the analytic derivations, independent direct-payoff audit, and regression suite rather than by Lean alone.
